\documentclass[10pt]{article}
\usepackage[margin=20mm]{geometry}
\usepackage{amsmath,amssymb,fontspec,unicode-math,xurl,hyperref}
\IfFontExistsTF{DejaVu Serif}{\setmainfont{DejaVu Serif}}{\setmainfont{Latin Modern Roman}}
\setmathfont{latinmodern-math.otf}
\setlength{\parindent}{0pt}\setlength{\parskip}{6pt}\sloppy
\begin{document}
\section*{Scientific verification report: I25-40119}
\section*{1. Identity and exact reviewed version}
Paper: Position: Generative AI Regulation Can Learn from Social Media Regulation. Author: Ruth Elisabeth Appel. ICML 2025, PMLR 267:81035-81049. Assigned display number: 0019.\par
The reviewed version is the 15-page final proceedings PDF, SHA-256 a8d0ff2846358b38cfa6922b2353fc40feac2f360fab4e98040474bea9c09f08, downloaded at 2026-10-24T04:36:28.808563+00:00. References to the paper use 1-based PDF pages, matching its printed pages. The earlier 11-page arXiv manuscript is distinguished throughout. [Appel-PMLR (2025)]\par
Review date: 24 October 2026 (UTC).\par
\section*{2. Summary and supported judgment}
Supported outcome: No Issue Found within the inspected scope. The review did not establish a substantive defect that invalidates the paper’s central position or warrants a numbered Problem entry. Selected empirical citations support their bounded claims, and historical legal texts support a qualified reading of the legal examples.\par
This finding is limited. The paper is a policy argument about learning from experience, with particular attention to content moderation and regulation broadly understood to include self-regulation. It does not demonstrate the effectiveness of a finished regulatory package. Several qualifications below matter for interpreting its evidence: model scoring is distinct from generated behavior; research-access mechanisms differ across laws; and medium-level feature comparisons are not universal technological classifications.\par
All 15 pages were read, including the bibliography. Table 1 on p.3 and Figure 1 on p.5 were visually inspected. Selected empirical sources, original historical legal texts, accessible version records, and public-material links were checked. Not every citation, legal jurisdiction, or policy outcome was independently audited. No new author experiment, theorem, dataset, or executable implementation is presented for this review to reproduce.\par
\section*{3. Argument reconstruction and quantifier check}
The central inference is that shared content-governance problems make social-media experience informative for generative-AI policy. That inference can support candidate practices and questions without implying that a practice will have identical effects in the two settings. The distinction matters: shared features alone do not identify an intervention’s net benefit when incentives, user populations, failure modes, and implementation costs differ.\par
The paper limits its comparison to common applications on p.2, discusses differences on pp.3-4, acknowledges institutional shortcomings on pp.5-6, calls for effectiveness research on p.7, and distinguishes learning from copying on p.9. These qualifications directly answer several possible adverse readings. An imperfect analogy is not sufficient to refute the modest proposition that relevant lessons can be drawn. [Appel-PMLR (2025), pp.2-9]\par
The four recommendations are policy directions. Figure 1 maps them to strategies, possible measures, and social-media precedents; it is not a plot of measured effects. Table 1 is a conceptual classification, not an estimated prevalence table or an experimental comparison. No numerical significance or causal-effect estimate was inferred from either visual.\par
The review therefore does not convert acknowledged future research into an alleged missing experiment. The broad comparative phrase that this is the best way to avoid mistakes is not established by a comparative policy evaluation; it is most defensibly read as advocacy within a paper that expressly permits other analogies and first-principles analysis.\par
\section*{4. Selected empirical-source checks}
\subsection*{4.1 Opinion distributions and political leaning}
The primary Santurkar et al. publication record describes comparisons between language-model responses and opinion distributions from 60 US demographic groups. It reports substantial misalignment and left-leaning tendencies in some human-feedback-tuned models. This supports the reviewed paper’s bounded discussion of opinion mismatch. It does not establish a context-independent orientation for all models. [Santurkar et al. (2023)]\par
The paper’s p.5 acknowledgment that prompts and question order affect bias measurement is relevant counterevidence against an interpretation that it ignores elicitation dependence. This verification checked the primary publication record and abstract; it did not recompute OpinionQA measurements.\par
\subsection*{4.2 Influence on opinions and decisions}
Fisher et al.’s 2024 version describes randomized interactive tasks involving liberal-biased, conservative-biased, and control chatbot conditions. The inspected methods include instruction-induced partisan behavior and tasks involving opinion formation and budget allocation. Those experiments support the possibility that biased chatbot behavior can affect choices in the tested setting. [Fisher et al. (2024), v1, section3.2]\par
This evidence does not quantify changes in actual election outcomes, establish identical effects for default deployed models, or prove that a proposed regulation removes the effect. The reviewed paper’s possibility claim does not require those stronger conclusions. The 2024 version was inspected, rather than substituting the source’s current 2026 revision.\par
Matz et al.’s original article was inspected; it studies personalized generated messages and outcomes involving attitudes and intended behavior. That evidence fits the bounded persuasion claim; it should not be read as direct observation of lasting real-world voting behavior. One Nature-page retrieval failed, but a separate retrieval obtained the full primary article. No reported effect size was independently recalculated. [Matz et al. (2024)]\par
\subsection*{4.3 CrowS-Pairs: measurement and model-class qualification}
Nangia et al. evaluate supplied stereotypical and less-stereotypical sentence pairs using masked-language-model scores. Their methods use conditional pseudo-log-likelihood, and the evaluated models are BERT-base, RoBERTa-large, and ALBERT-XXL-v2. [Nangia et al. (2020), sections3-4, PDF pp.4-5 / printed pp.1956-1957]\par
This is evidence of biases in the tested language models’ sentence-scoring behavior. It is not a direct experiment measuring freely generated depictions from a contemporary conversational or image-generation product. The citation on p.4 is best framed as earlier evidence about language-model biases, with modern deployment-specific claims supported separately.\par
The charitable reading is that the citation supplies background evidence for the broader bias concern. Other cited studies support that concern, so this source-type qualification does not establish that the policy premise is false or that its main recommendations fail. It is recorded as a scope clarification, without inventing a corrected empirical effect.\par
\section*{5. Historical legal-source checks}
\subsection*{5.1 Access mechanisms must be distinguished}
The historical AI Act text permits specified general-purpose-model compliance and systemic-risk evaluations, including independent experts appointed to act for the Commission. Article 92 permits access requests through APIs or other appropriate technical means. This supports a narrow reading of the p.8 reference to researcher access, but does not create a general access entitlement for unaffiliated researchers’ own projects. [AI Act signed text (2024), Article92, printed pp.342-343 / PDF pp.343-344]\par
The DSA has a different route: Article40 provides conditional access for vetted researchers, through the relevant Digital Services Coordinator, for systemic-risk and mitigation research involving covered very large platforms/search engines. Article37 requires annual independent audits of those covered services. These are bounded mechanisms, not duties applying identically to every service and every research request. [DSA signed text (2022), Articles37 and40]\par
Legal example | Verified interpretation | Limit on transfer\par
AI Act Article92 | Regulator-directed evaluation, potentially involving appointed independent experts | Not a general independent-research application right\par
DSA Article40 | Conditional researcher access for designated purposes and covered providers | Eligibility, procedure, purpose, and provider scope matter\par
DSA Article37 | Independent audits for covered very large services | Not an undifferentiated audit requirement for all social media\par
The paper does not expressly say these access mechanisms are legally identical. Its recommendation to improve access remains intelligible under the narrower reading. The appropriate outcome is legal-scope qualification, not a fabricated contradiction.\par
\subsection*{5.2 User-base factors and historical timing}
AI Act Annex XIII(f)-(g) includes market reach and registered-user factors, including a presumption based on at least 10,000 registered EU business users. The p.8 reference to user-base size therefore has support; it is not evidence of accidental substitution of a DSA threshold. [AI Act signed text (2024), AnnexXIII, PDF p.419]\par
The verified legal sources are signed Council/Parliament texts dated 13 June 2024 and 19 October 2022. Direct EUR-Lex retrieval encountered a JavaScript verification page; the official Council copies were accessible. Their precise provenance is retained rather than presenting them as a retrieved current consolidated-law edition.\par
This review does not certify current compliance, resolve the application of Section230 to generated speech, or equate historical enactment with the date every obligation became applicable. The paper itself presents liability questions as uncertain. Subsequent legal or company-policy changes are not used to manufacture errors in its historical account.\par
\section*{6. Technical comparison and policy-effect boundaries}
The comparison between deep generative models and traditional recommender systems should not be interpreted as a division between deep learning and all recommendation systems. A primary 2016 YouTube paper already describes deep candidate-generation and ranking models. The distinction between a probabilistic model and a stochastic output-selection procedure also cuts across content types. [Covington et al. (2016)]\par
For a simple analytic illustration, a fixed distribution assigning probabilities 0.7 and 0.3 to two options can be used with deterministic argmax or with categorical sampling. The options could be text tokens or recommended items. Thus probabilistic representation and randomized selection are separate properties. This is a conceptual illustration, not a simulation of a deployed platform.\par
The paper’s p.2 caveat about common applications, its comparison of differences in degree, and its limited content-moderation focus prevent a single atypical application from automatically refuting the argument. More precise descriptions would identify the product class, date, sampling and ranking pipeline, and stage of moderation. The present review did not establish a central failure caused by the simplified binary feature labels.\par
Similarly, transparency, democratic participation, choice, staffing, and global expertise are candidate safeguards, not interchangeable outcomes. An implementation assessment would need to specify the harm, affected population, intervention, comparison, costs, privacy constraints, and success criteria. The paper explicitly requests research on effectiveness, so those unresolved questions are part of its stated agenda.\par
\section*{7. Version history, materials, and known-issue status}
The original PDF matches an immutable PMLR copy at commit a3e4112da70e3f0623daeb3eeea06e2b34ed5e2f. The arXiv record 2412.11335 lists v1 submitted on 15 December 2024 at 23:00:29 UTC, associated with the NeurIPS2024 Regulatable ML workshop. It is an earlier 11-page manuscript. The saved earlier PDF has SHA-256 8f57b140ed7a74f257211f954d17bcca1ef6e5f3fd819a1c4bafb2cfe3ca1c7b. [Appel-arXiv (2024)]\par
The author website’s full-text link resolves to that earlier arXiv record. It should not replace the proceedings source or its pagination. An exhaustive semantic or visual comparison of versions was not performed.\par
No author erratum or later corrected edition was identified in the inspected PMLR, arXiv, or author publication records. This is not a universal absence claim. The official OpenReview forum and a related workshop record returned verification pages, leaving their discussions, attachments, and revision histories unverified. No paper-specific code/data package was found in the inspected links; for this policy position paper, that absence alone is not a reproducibility defect.\par
\section*{8. Conclusion table}
Item | Judgment | Verification or repair status | Adoption | Remaining work\par
Core learning-from-experience position | No Issue Found within inspected scope | Argument and stated limitations checked; no repair established as necessary | Not applicable | Human assessment and specific implementation evaluation\par
Selected opinion/persuasion evidence | No Issue Found within selected source checks | Source support checked; original estimates not reanalyzed | Not applicable | Broader citation and external-validity audit\par
CrowS-Pairs use | Scope clarification | Model class and scoring metric checked | Not requested | Use deployment-specific evidence for deployed-output claims\par
Legal access and audit examples | Scope clarification | Historical provisions checked with covered actors/purposes | Not requested | Jurisdiction-specific implementation; current compliance outside scope\par
Effectiveness of policy proposals | Unresolved future work acknowledged by author | No efficacy estimate or repair claimed | Not applicable | Empirical testing with defined outcomes and comparators\par
No Major or Minor finding is asserted. No author adoption or human scientific approval is recorded.\par
\section*{9. Consolidated references}
[Appel-PMLR (2025)] Appel, R.E. PMLR 267:81035-81049. https://proceedings.mlr.press/v267/appel25a.html . Immutable PDF: https://raw.githubusercontent.com/mlresearch/v267/a3e4112da70e3f0623daeb3eeea06e2b34ed5e2f/assets/appel25a/appel25a.pdf .\par
[Santurkar et al. (2023)] Whose Opinions Do Language Models Reflect? PMLR 202:29971-30004. https://proceedings.mlr.press/v202/santurkar23a.html .\par
[Fisher et al. (2024)] Biased AI can Influence Political Decision-Making. arXiv:2410.06415v1. https://arxiv.org/abs/2410.06415v1 and https://arxiv.org/html/2410.06415v1 .\par
[Matz et al. (2024)] The potential of generative AI for personalized persuasion at scale. Scientific Reports 14:4692. https://doi.org/10.1038/s41598-024-53755-0 .\par
[Nangia et al. (2020)] CrowS-Pairs: A Challenge Dataset for Measuring Social Biases in Masked Language Models. EMNLP, pp.1953-1967. https://aclanthology.org/2020.emnlp-main.154/ .\par
[AI Act signed text (2024)] European Parliament and Council, PE-CONS24/1/24 REV1, 13 June2024. https://data.consilium.europa.eu/doc/document/PE-24-2024-REV-1/EN/pdf . Targeted Articles92 and AnnexXIII; not a current-law compliance opinion.\par
[DSA signed text (2022)] European Parliament and Council, PE-CONS30/1/22 REV1, 19 October2022. https://data.consilium.europa.eu/doc/document/PE-30-2022-REV-1/en/pdf . Article37 printed p.212, Article40 printed pp.219-225.\par
[Covington et al. (2016)] Deep Neural Networks for YouTube Recommendations. RecSys2016. https://research.google/pubs/deep-neural-networks-for-youtube-recommendations/ .\par
[Appel-arXiv (2024)] arXiv:2412.11335v1. https://arxiv.org/abs/2412.11335 . Author listing: https://ruthappel.github.io/academic-writing/ . Inaccessible forum: https://openreview.net/forum?id=fRk0nKLKrJ .\par
\section*{Appendix A. Actual procedures, source manifests, and failures}
The original download, SHA-256, byte count, UTC timestamp, and successful pdftotext/-layout commands are preserved in sources/0019/source\_download.json. Table1 and Figure1 were rendered using pdftoppm -f PAGE -l PAGE -singlefile -r120 -png for pages3 and5; both commands exited0. The images and command log are preserved in sources/0019.\par
The independent retrieval script sources/0019/retrieve\_support.py ran with Python -B and exited0. It downloaded five supporting PDFs, checked their PDF headers, calculated hashes, and extracted text successfully: CrowS-Pairs, the 2016 YouTube article, the two signed Council legislative texts, and the earlier arXiv manuscript. The full URL/hash/size/command records are in sources/0019/supporting\_downloads.json. These source retrievals are not empirical replications.\par
Access failures: direct EUR-Lex pages required JavaScript verification; official Council copies supplied the historical text. OpenReview forums remained inaccessible. One Nature retrieval redirected to an inaccessible authorization endpoint, while a separate retrieval obtained the primary article. The Jakesch et al. ACM DOI returned403; full-text verification of that citation is not claimed. Verification challenges were not bypassed.\par
An initial coverage.json write failed because the numbered verification directory was absent. Directory-state inspection confirmed this, after which it was created and the artifact writes completed. No author code, model inference, statistical reanalysis, or package installation occurred. The analytic probability example in section6 is a hand-derived illustration, not a reported software experiment.\par
\section*{Supporting checks}
The independent check read the full paper, inspected its visuals, checked primary evidence and legal excerpts, and synthesized this report. The visible app selection was not independently inspected.\par
Remaining limits include incomplete audit of all cited studies, no evaluation of actual policy effectiveness, unresolved OpenReview material, and no exhaustive legal analysis across jurisdictions. These limits do not amount to proof of a substantive defect.\par
Human scientific decisions and publication status remain pending.\par
\end{document}
